\documentclass[11pt,a4paper]{article}

\usepackage[T1]{fontenc}
\usepackage[utf8]{inputenc}
\usepackage{lmodern}
\usepackage{amssymb,amsmath,amsthm,mathtools}
\usepackage[margin=2.5cm]{geometry}
\usepackage{parskip}
\usepackage[hidelinks]{hyperref}
\usepackage{microtype}

% Theorem environments (arXiv / journal style)
\theoremstyle{plain}
\newtheorem{theorem}{Theorem}[section]
\newtheorem{proposition}[theorem]{Proposition}
\newtheorem{lemma}[theorem]{Lemma}
\newtheorem{corollary}[theorem]{Corollary}

\theoremstyle{definition}
\newtheorem{definition}[theorem]{Definition}
\newtheorem{example}[theorem]{Example}
\newtheorem{remark}[theorem]{Remark}

% Notation
\newcommand{\Ext}{\mathrm{Ext}}
\newcommand{\Shape}{\mathrm{Shape}}
\newcommand{\Pos}{\mathrm{Pos}}
\newcommand{\Type}{\mathbf{Type}}
\newcommand{\id}{\mathrm{id}}
\newcommand{\rootp}{\mathrm{o}}

\title{From Container Comonads Back to Directed Containers}
\author{Extracted from a machine-checked development}
\date{}

\begin{document}
\maketitle

\begin{abstract}
We prove the converse of the Ahman--Chapman--Uustalu correspondence between
directed containers and comonads on container extensions. Starting from a
comonad structure on the extension of a container, we recover root, sub-shape,
and shift operations, and show that the five directed-container axioms follow
from the three comonad laws. A key step is that the right counit law forces the
outer shape map of the comultiplication to be the identity. The proofs of all
stated results have been machine-checked; we record only the mathematical
content.
\end{abstract}

\tableofcontents
\bigskip

%======================================================================
\section{Introduction}
%======================================================================

Ahman, Chapman, and Uustalu showed that a container extension carries a
comonad structure if and only if the container is \emph{directed}: it is
equipped with a root, a sub-shape operation, and a shift satisfying five
equational axioms D1--D5. The forward direction is standard: from directed
data one builds counit and comultiplication on the extension. This note
treats the converse.

By the representation theorem for containers, natural transformations of
extensions are container morphisms. Consequently a comonad on $\Ext\,C$ is
determined by combinatorial data: a root choice from the counit, and a triple
of maps from the comultiplication. The right counit law forces the outer shape
component of the comultiplication to be the identity. Once that collapses, the
remaining data are precisely root, sub, and shift, and the three comonad laws
become D1--D5---the same partition as in the forward direction, read in
reverse.

We work in a single universe of types. Proofs are indicated by the phrase
``by lean,'' meaning they have been checked by a proof assistant rather than
written out in prose.

\paragraph{What this note claims.}
A comonad on a container extension is, up to the transport forced by a
propositional identity of shapes, exactly a directed container. More
precisely: from counit and comultiplication data satisfying the three comonad
laws one recovers operations $(\rootp,\downarrow,\oplus)$ satisfying D1--D5;
packaging a directed container as a comonad and recovering returns the
original directed container; and the recovered structure induces the same
counit and comultiplication as those one started from.

%======================================================================
\section{Containers, extensions, and generic elements}
%======================================================================

\begin{definition}[Container and extension]
A \textbf{container} $C$ consists of a type $\Shape(C)$ of shapes and, for
each shape~$s$, a type $\Pos(C)(s)$ of positions. The \textbf{extension} at a
type~$X$ is
\begin{equation}
  \Ext\,C\,X \;=\; \sum_{s \,:\, \Shape(C)}\;
    \bigl(\Pos(C)(s) \to X\bigr).
\end{equation}
We write $\langle s,\, v\rangle$ for a typical element.
\end{definition}

\begin{definition}[Functorial action]
For $f : X \to Y$,
\begin{equation}
  \Ext\,C\,f\,\bigl\langle s,\, v\bigr\rangle
  \;=\;
  \bigl\langle s,\; f \circ v\bigr\rangle.
\end{equation}
\end{definition}

The following two lemmas package elementary structure of dependent pairs that
will be used to read laws off equalities of extensions.

\begin{lemma}[Injectivity of extension pairs]\label{lem:mk-inj}
If $\langle s_1,\, v_1\rangle = \langle s_2,\, v_2\rangle$ in $\Ext\,C\,X$,
then there exists $h : s_1 = s_2$ such that
$v_1(q) = v_2(h \mathbin{\triangleright} q)$ for all positions~$q$ of~$s_1$.
\end{lemma}

\begin{proof}
By lean.
\end{proof}

\begin{lemma}[Normalisation along a shape equality]\label{lem:norm-aux}
Let $t = s$ be a shape equality, let $\mathrm{sub}$ assign shapes to positions
of~$t$, and let $\mathrm{val}$ map nested positions back to positions of~$s$.
Then for any labelling~$v$ of~$s$,
\begin{equation}
\begin{aligned}
  &\bigl\langle t,\;
    \lambda p.\,
      \bigl\langle \mathrm{sub}\,p,\;
        \lambda q.\, v(\mathrm{val}\,p\,q)\bigr\rangle
  \bigr\rangle \\
  &\qquad=\;
  \bigl\langle s,\;
    \lambda p.\,
      \bigl\langle \mathrm{sub}\,(h^{-1} \mathbin{\triangleright} p),\;
        \lambda q.\,
          v\bigl(\mathrm{val}\,(h^{-1} \mathbin{\triangleright} p)\,q\bigr)
      \bigr\rangle
  \bigr\rangle,
\end{aligned}
\end{equation}
where $h$ is the given equality $t = s$.
\end{lemma}

\begin{proof}
By lean.
\end{proof}

%======================================================================
\section{Combinatorial data of a container comonad}
%======================================================================

By the representation theorem, a natural transformation
$\Ext\,C \Longrightarrow \Ext\,D$ is a container morphism. Applying this to
counit and comultiplication yields the following packaging.

\begin{definition}[Container comonad data]\label{def:container-comonad}
A \textbf{container comonad} is a container equipped with
\begin{enumerate}
  \item a root choice $e : (s : S) \to P(s)$;
  \item an outer shape map $c_S : S \to S$;
  \item an inner shape assignment
        $c_{\downarrow} : (s : S) \to P(c_S(s)) \to S$;
  \item a value map
        $c_{\oplus} : (s : S) \to (p : P(c_S(s))) \to
          P(c_{\downarrow}(s,p)) \to P(s)$.
\end{enumerate}
\end{definition}

\begin{definition}[Induced counit and comultiplication]\label{def:counit-comult}
Given container comonad data~$M$, define
\begin{align}
  \varepsilon\langle s,\, v\rangle
  &\;=\;
  v\bigl(e(s)\bigr), \\[4pt]
  \delta\langle s,\, v\rangle
  &\;=\;
  \bigl\langle
    c_S(s),\;
    \lambda p.\,
      \bigl\langle
        c_{\downarrow}(s,p),\;
        \lambda q.\, v\bigl(c_{\oplus}(s,p,q)\bigr)
      \bigr\rangle
  \bigr\rangle.
\end{align}
\end{definition}

\begin{definition}[Comonad laws]\label{def:is-comonad}
The data~$M$ \textbf{satisfy the comonad laws} when, for every type~$X$,
\begin{align}
  \varepsilon \circ \delta &\;=\; \id &&\text{(left counit),} \\
  (\Ext\,C\,\varepsilon) \circ \delta &\;=\; \id &&\text{(right counit),} \\
  (\Ext\,C\,\delta) \circ \delta &\;=\; \delta \circ \delta
    &&\text{(coassociativity).}
\end{align}
\end{definition}

%======================================================================
\section{Collapse of the outer shape}
%======================================================================

\begin{theorem}[Outer shape is the identity]\label{thm:cshape-eq}
If $M$ satisfies the comonad laws, then for every shape~$s$,
\begin{equation}
  c_S(s) \;=\; s.
\end{equation}
\end{theorem}

\begin{proof}
By lean.
\end{proof}

\begin{remark}[Idea of the argument]
Probe the right counit law $(\Ext\,C\,\varepsilon) \circ \delta = \id$ at the
generic element $\langle s,\, \id\rangle \in \Ext\,C\,(P(s))$. The left-hand
side has outer shape $c_S(s)$; the right-hand side has outer shape~$s$. Reading
off the shape component yields the claim.
\end{remark}

%======================================================================
\section{Recovered directed-container operations}
%======================================================================

Using Theorem~\ref{thm:cshape-eq} to transport positions into the $c_S(s)$
fibre, we read root, sub-shape, and shift off the combinatorial data.

\begin{definition}[Recovered operations]\label{def:recovered}
Assume $M$ satisfies the comonad laws. Define
\begin{align}
  \rootp(s) &\;=\; e(s), \\[4pt]
  s \downarrow p
  &\;=\;
  c_{\downarrow}\bigl(s,\; (c_S(s)=s)^{-1} \mathbin{\triangleright} p\bigr),
  \\[4pt]
  p \oplus q
  &\;=\;
  c_{\oplus}\bigl(s,\;
    (c_S(s)=s)^{-1} \mathbin{\triangleright} p,\; q\bigr).
\end{align}
\end{definition}

\begin{proposition}[Normalised comultiplication]\label{prop:comult-norm}
For every shape~$s$ and labelling~$v$,
\begin{equation}
  \delta\langle s,\, v\rangle
  \;=\;
  \bigl\langle s,\;
    \lambda p.\,
      \bigl\langle s \downarrow p,\;
        \lambda q.\, v(p \oplus q)\bigr\rangle
  \bigr\rangle.
\end{equation}
\end{proposition}

\begin{proof}
By lean.
\end{proof}

Thus, once the outer shape collapses, comultiplication is already in the
standard directed-container form, up to the transports inherent in
Definition~\ref{def:recovered}.

%======================================================================
\section{Directed-container axioms from comonad laws}
%======================================================================

Each directed-container law is obtained by probing a comonad law at a generic
element and normalising the comultiplication. The partition matches the
forward direction run in reverse: left counit yields D1 and D2; right counit
yields D3; coassociativity yields D4 and D5.

\begin{theorem}[D1 from left counit]\label{thm:d1}
For every shape~$s$,
\begin{equation}
  s \downarrow \rootp(s) \;=\; s.
\end{equation}
\end{theorem}

\begin{proof}
By lean.
\end{proof}

\begin{theorem}[D2 from left counit]\label{thm:d2}
For every shape~$s$ and every position $q$ of $s \downarrow \rootp(s)$,
\begin{equation}
  \rootp(s) \oplus q \;=\; (s \downarrow \rootp(s) = s) \mathbin{\triangleright} q.
\end{equation}
\end{theorem}

\begin{proof}
By lean.
\end{proof}

\begin{theorem}[D3 from right counit]\label{thm:d3}
For every shape~$s$ and position~$p$ of~$s$,
\begin{equation}
  p \oplus \rootp\bigl(s \downarrow p\bigr) \;=\; p.
\end{equation}
\end{theorem}

\begin{proof}
By lean.
\end{proof}

\begin{theorem}[D4 from coassociativity]\label{thm:d4}
For every shape~$s$, position~$p$ of~$s$, and position~$q$ of
$s \downarrow p$,
\begin{equation}
  s \downarrow (p \oplus q) \;=\; \bigl(s \downarrow p\bigr) \downarrow q.
\end{equation}
\end{theorem}

\begin{proof}
By lean.
\end{proof}

\begin{theorem}[D5 from coassociativity]\label{thm:d5}
For positions as in D4 and a further nested position~$q'$,
\begin{equation}
  (p \oplus q) \oplus
    \bigl((s \downarrow (p \oplus q)
      = (s \downarrow p) \downarrow q)^{-1}
      \mathbin{\triangleright} q'\bigr)
  \;=\;
  p \oplus (q \oplus q').
\end{equation}
\end{theorem}

\begin{proof}
By lean.
\end{proof}

\begin{definition}[Recovery of a directed container]\label{def:to-directed}
Given container comonad data satisfying the comonad laws, the
\textbf{recovered directed container} has the same underlying container, root
$\rootp$, sub~$\downarrow$, and shift~$\oplus$ as in
Definition~\ref{def:recovered}, with laws D1--D5 given by
Theorems~\ref{thm:d1}--\ref{thm:d5}.
\end{definition}

\begin{theorem}[Recovery preserves the counit]\label{thm:counit-same}
The directed container of Definition~\ref{def:to-directed} induces the same
counit as the original container comonad.
\end{theorem}

\begin{proof}
By lean.
\end{proof}

\begin{theorem}[Recovery preserves the comultiplication]\label{thm:comult-same}
The directed container of Definition~\ref{def:to-directed} induces the same
comultiplication as the original container comonad.
\end{theorem}

\begin{proof}
By lean.
\end{proof}

Together, Theorems~\ref{thm:counit-same} and~\ref{thm:comult-same} say that
applying the forward construction to the recovered directed container returns
the original comonad.

%======================================================================
\section{The other half of the equivalence}
%======================================================================

Conversely, every directed container packages as container comonad data with
outer shape map the identity.

\begin{definition}[Packaging]\label{def:packaging}
Let $D$ be a directed container. Define container comonad data by
$e = \rootp$, $c_S = \id$, $c_{\downarrow} = \downarrow$, and
$c_{\oplus} = \oplus$.
\end{definition}

\begin{proposition}[Packaged data form a comonad]\label{prop:packaged-is-comonad}
The data of Definition~\ref{def:packaging} satisfy the three comonad laws.
\end{proposition}

\begin{proof}
By lean.
\end{proof}

\begin{theorem}[Round-trip on directed containers]\label{thm:roundtrip}
Packaging a directed container as a container comonad and recovering yields
the original directed container.
\end{theorem}

\begin{proof}
By lean.
\end{proof}

\begin{remark}
Because $c_S = \id$ definitionally in the packaging, every transport along
$c_S(s) = s$ is the identity, and the recovered operations agree with the
original ones without residual casts.
\end{remark}

%======================================================================
\section{Discussion}
%======================================================================

\paragraph{What has been established.}
Container comonads and directed containers determine each other. The linchpin
is Theorem~\ref{thm:cshape-eq}: comultiplication cannot change the outer
shape. After that collapse, comonad laws and directed-container axioms are
two presentations of the same equational content, partitioned identically in
both directions.

\paragraph{Scope.}
The development works with comonad data already presented in combinatorial
form (as guaranteed by the representation theorem). It does not re-prove that
every natural transformation arises from a container morphism; that is
assumed as background. Dependent transports along propositional equalities of
shapes are intrinsic to the converse and mirror those of the forward theory.

\paragraph{Machine checking.}
Every definition and theorem statement above corresponds to a declaration in a
machine-checked formalisation. The phrase ``by lean'' indicates that the proof
has been verified by a proof assistant.

%======================================================================
\section*{Acknowledgements}
%======================================================================

The mathematical content follows Ahman, Chapman, and Uustalu's characterisation
of container comonads. The presentation here is a prose rendering of a formal
development intended for human review of the theorem statements.

\begin{thebibliography}{9}

\bibitem{ahman2014}
D.~Ahman, J.~Chapman, and T.~Uustalu.
\newblock When is a container a comonad?
\newblock \emph{Logical Methods in Computer Science}, 10(3), 2014.

\bibitem{abbott2005}
M.~Abbott, T.~Altenkirch, and N.~Ghani.
\newblock Containers: constructing strictly positive types.
\newblock \emph{Theoretical Computer Science}, 342(1):3--27, 2005.

\end{thebibliography}

\end{document}
